% MAS1004 Assignment 1: long report.
%
% You fill this in with your agent, in your repository. AGENTS.md tells the
% agent how. Red text in the PDF is what still has to be replaced.
%
% - Pictures go in figures/. Copy them from results/, which is not in git.
% - Outputs go in outputs/. run.py saves what it prints in
%   results/<tag>_output.txt, and PLAN.md saves the other outputs in results/
%   with `| tee results/...`. Copy them here. \outputfile{outputs/check.txt}
%   prints a file exactly as it is. Every number in this report has to come
%   from such an output.
%
% To make the PDF: compress this folder into a zip, and on Overleaf choose
% New Project, Upload Project. Set Menu, Compiler to XeLaTeX, and Recompile.

\documentclass[11pt]{article}
\usepackage{mas1004}

\reporttitle{Assignment 1: Long report}
\studentname{Your name}
\studentid{Your student ID}

\begin{document}
\makeheader

% ---------------------------------------------------------------------------
\section*{Problem 1. Categories and images}

\subsection*{Why these categories}
\fillin{Why these categories, and why you care about them.}

\subsection*{What I expected the model to use}
\fillin{What you thought, before training, the model would use to tell them
apart: shape, colour, texture, background, or something else.}

\subsection*{Images per category}
\fillin{How many images you got for each category.}

% ---------------------------------------------------------------------------
\section*{Problem 2. Preparing an image}

\sidebyside{figures/prepare_original.jpg}{\fillin{the image as downloaded}}
           {figures/prepare_result.png}{\fillin{what prepare\_image made of it}}

\fillin{What was cut off, and an image where the crop cut off part of the
thing you care about.}

% ---------------------------------------------------------------------------
\section*{Problem 3. First training runs}

\begin{center}\footnotesize
\begin{tabular}{llllrrrr}
\toprule
tag & start & trained & epochs & lr & trainable parameters & train acc. & test acc. \\
\midrule
run & ImageNet & all layers & 10 & 0.0001 & 11,178,564 & 99.8\% & 88.3\% \\
scratch & random & all layers & 10 & 0.0001 & 11,178,564 & 99.6\% & 65.0\% \\
frozen & ImageNet & last layer & 10 & 0.001 & 2,052 & 94.6\% & 88.3\% \\
custom & ImageNet & all layers & 15 & 2e-06 & 11,178,564 & 72.5\% & 60.8\% \\
\bottomrule
\end{tabular}
\end{center}

\picturefile{figures/run_curves.png}
\picturefile{figures/scratch_curves.png}
\picturefile{figures/frozen_curves.png}
\picturefile{figures/custom_curves.png}

\fillin{Did the loss go down? Did it keep going down, or flatten out?}

\fillin{The gap between train and test accuracy, and what it means.}

\begin{center}\picturefile[0.5\linewidth]{figures/run_confusion.png}\end{center}
\begin{center}\picturefile[0.5\linewidth]{figures/scratch_confusion.png}\end{center}
\begin{center}\picturefile[0.5\linewidth]{figures/frozen_confusion.png}\end{center}
\begin{center}\picturefile[0.5\linewidth]{figures/custom_confusion.png}\end{center}

\fillin{Which two categories it mixes up most, and whether that surprised you.}

\fillin{How far apart the test accuracies are, starting from ImageNet and from
random numbers, and why you think that is.}

\fillin{How close training only the last layer gets to training all of them.}

\subsection*{Output of the runs}
\outputfile{outputs/run_output.txt}
\outputfile{outputs/scratch_output.txt}
\outputfile{outputs/frozen_output.txt}
\outputfile{outputs/custom_output.txt}

% ---------------------------------------------------------------------------
\section*{Problem 4. Cleaning}

\subsection*{My rule}
\begin{description}
\item[\texttt{tooth brush}] 솔이 몸체보다 두껍지 않고 몸체보다 훨씬 작아야 한다. 솔의 반대편은 평면이어야 하며, 솔이 동그란 원통 모양인 것은 속하지 않는다.
\item[\texttt{hair brush}] 솔과 몸체의 비율이 1:1 정도로 비슷하다. 솔이 한쪽 면에만 있는 종류와 원통 모양으로 동그랗게 생긴 종류가 있으며, 솔 사이에 간격이 있다.
\item[\texttt{cleaning brush}] 손잡이가 있고, 솔보다 몸체의 길이가 길다. 손잡이에 몸체와 다른 색의 고무패킹이 있는 경우가 있으며, 솔과 몸체가 이어지는 부분은 약 130도 정도로 휘어져 있다.
\item[\texttt{shoe brush}] 몸체 전체에 솔이 있는 경우가 많다. 나무 재질이 연상되는 갈색 몸체를 가지며 솔은 일반적으로 검정색이다.
\end{description}

\subsection*{Three kinds of junk}
\fillin{Three kinds of junk, with one picture of each.}

% For example: \picturefile[0.3\linewidth]{figures/junk_1.jpg}

\subsection*{What I removed}
\outputfile{outputs/clean_count.txt}

The check found 35 pairs of near copies. I moved 30 right-side images to
\texttt{data/removed/}, keeping one representative from each connected group.
I kept \texttt{shoe\_brush/0049.jpg} because its left-side copy,
\texttt{cleaning\_brush/0091.jpg}, was removed as a mislabeled shoe brush.

Among the first 20 suspects, I removed 5 \texttt{cleaning\_brush}, 1
\texttt{hair\_brush}, 5 \texttt{shoe\_brush}, and 0 \texttt{tooth\_brush}
images. Reviewing the full Step 3 sheets also revealed six more cleaning-brush
images outside the 32 shown on the suspect sheet:
\texttt{0031.jpg}, \texttt{0112.jpg}, \texttt{0128.jpg}, \texttt{0134.jpg},
\texttt{0137.jpg}, and \texttt{0150.jpg}.

\subsection*{Before and after cleaning}
\outputfile{outputs/clean_output.txt}
\picturefile{figures/clean_curves.png}
\begin{center}\picturefile[0.5\linewidth]{figures/clean_confusion.png}\end{center}
\picturefile{figures/clean_worst.png}
\outputfile{outputs/compare.txt}

The downloaded-test accuracy changed from 88.3\% (\texttt{run}) to 93.6\%
(\texttt{clean}). On the same 25 new images, \texttt{run} got 19 correct
(76.0\%) and \texttt{clean} got 21 correct (84.0\%). The new-image comparison
is fairer because both models are measured on exactly the same images, while
cleaning changed the downloaded test set.

% ---------------------------------------------------------------------------
\section*{Problem 5. Images from a new source}

\fillin{Where your new images came from, and why you are sure none of them are
among your downloads.}

\outputfile{outputs/overlap.txt}
\outputfile{outputs/check.txt}

The cleaned model scored 93.6\% on the downloaded test set and 84.0\% on the
new images (21/25), as printed by the training and check outputs.

\fillin{Five mistakes it was confident about, with the pictures. For each one,
what in the image you think pushed it the wrong way.}

\fillin{What is different about your new images: background, lighting, angle,
distance, what else is in the frame.}

% ---------------------------------------------------------------------------
\section*{Problem 6. The demo}

Address: \fillin{your GitHub Pages address}

\fillin{What happens when you show it something that is none of your
categories.}

% ---------------------------------------------------------------------------
\section*{Problem 7. What the model looks at}

% Copy this subsection once for every guess.
\subsection*{Guess 1}

\fillin{The guess, and what in your images or your model's mistakes made you
think of it.}

\fillin{The change you chose, and why it tests the guess.}

My request to the agent, word for word:
\begin{quote}
\fillin{the request}
\end{quote}

% One \sidebyside for every image you changed.
\sidebyside{figures/guess1_original_1.jpg}{\fillin{original: the model's answer}}
           {figures/guess1_changed_1.jpg}{\fillin{changed: the model's answer}}

\fillin{What the results say about the guess.}

\subsection*{What I now think the model uses}
\fillin{What you now think the model uses to tell your categories apart, and
how sure you are.}

% ---------------------------------------------------------------------------
\section*{Appendix. Output of check.py}
\outputfile{outputs/check.txt}

\end{document}
